% Source: https://artofproblemsolving.com/community/c3h2146379_pirates
% Stable identity: aops-topic-2146379-q01
Write $N=2015!$. The exchanges are numbered $2,3,\dots,2015$: in exchange $k$, Jack gives $\dfrac{1}{k}$ of his gold to Will if $k$ is even, and Will gives $\dfrac{1}{k}$ of his gold to Jack if $k$ is odd. All the amounts below are integers, because $N$ is divisible by every number up to $2015$.

\textbf{Claim.} After each odd-numbered exchange, both have $N$ again.

Suppose both have $N$ before exchange $2j$ (true at the start, $j=1$). In exchange $2j$, Jack gives $\dfrac {N}{2j}$, leaving
\[\text{Jack}=N\cdot\dfrac{2j-1}{2j},\qquad \text{Will}=N\cdot\dfrac{2j+1}{2j}.\]
In exchange $2j+1$, Will gives $\dfrac{1}{2j+1}$ of his gold, which is $\dfrac{N}{2j}$. Then Will has $N\cdot\dfrac{2j+1}{2j}-\dfrac {N}{2j}=N$ and Jack has $N\cdot\dfrac{2j-1}{2j}+\dfrac {N}{2j}=N$. This proves the claim by induction.

The last exchange, number $2015$, is odd-numbered, so at the end both have $N$, and Will has $\boxed{0}$ more pieces of gold than Jack.
